\documentclass[11pt,a4paper]{article}
\usepackage[margin=2.2cm,top=2.0cm,bottom=2.2cm]{geometry}
\usepackage{amsmath,amssymb}
\usepackage[T1]{fontenc}
\usepackage{mathptmx}
\usepackage{xcolor}
\usepackage{enumitem}
\usepackage{fancyhdr}
\usepackage{tabularx}
\usepackage{mdframed}
\usepackage{multicol}

\definecolor{navy}{HTML}{1F4E79}
\definecolor{rule}{HTML}{333333}
\definecolor{boxbg}{HTML}{F4F6F9}
\definecolor{keyred}{HTML}{B02020}

% \ANSWERStrue for the tutor key
\newif\ifANSWERS
\input{mode}

\pagestyle{fancy}\fancyhf{}
\renewcommand{\headrulewidth}{0pt}
\renewcommand{\footrulewidth}{0.4pt}
\fancyfoot[L]{\footnotesize\color{navy}SAT Math: Foundations to Mastery \ $\cdot$\ Signs and Checking\ifANSWERS\ $\cdot$\ \textbf{\color{keyred}ANSWER KEY}\fi}
\fancyfoot[R]{\footnotesize\color{navy}Page \thepage\ of \pageref{LastPage}}
\usepackage{lastpage}

\setlength{\parindent}{0pt}
\setlength{\parskip}{4pt}

\newcommand{\PartHead}[2]{%
  \vspace{6pt}
  {\color{navy}\rule{\textwidth}{0.8pt}}\\[2pt]
  {\large\bfseries\color{navy}#1}\hfill{\small\itshape #2}\\[-4pt]
  {\color{navy}\rule{\textwidth}{0.4pt}}
  \vspace{4pt}
}
\newcommand{\sitting}[1]{{\footnotesize\textsc{\color{navy}#1}}\par\vspace{2pt}}
\newcommand{\spr}{\\[3pt]\hspace*{6pt}Answer: \underline{\hspace{3.2cm}}\hfill{\footnotesize\itshape(write your answer)}}
\ifANSWERS
  \newcommand{\ans}[1]{\ {\color{keyred}\bfseries #1}}
  \newcommand{\note}[1]{\par{\footnotesize\color{keyred}#1}}
\else
  \newcommand{\ans}[1]{}
  \newcommand{\note}[1]{}
\fi

\begin{document}
\label{firstpage}

\begin{center}
{\LARGE\bfseries\color{navy}Signs and Checking}\\[4pt]
{\color{navy}SAT Math $\cdot$ Foundations to Mastery $\cdot$ Week of 27 September}
\ifANSWERS\\[4pt]{\color{keyred}\bfseries ANSWER KEY --- tutor copy}\fi
\end{center}
\vspace{2pt}
{\color{rule}\rule{\textwidth}{1pt}}
\vspace{6pt}

\ifANSWERS\else
Name: \underline{\hspace{6cm}} \hfill Due: \underline{\hspace{3.2cm}}
\vspace{6pt}
\fi

\begin{mdframed}[backgroundcolor=boxbg,linecolor=navy,linewidth=0.6pt,innertopmargin=6pt,innerbottommargin=8pt]
{\bfseries\color{navy}How to use this sheet}\par\vspace{2pt}
This is four short sittings, not one long one. Four half-hours across the week will
do more for you than one long Saturday, and you are far more likely to finish it.

\begin{itemize}[leftmargin=14pt,itemsep=1pt,topsep=3pt]
  \item \textbf{Part A is no calculator.} None. It is 25 minutes and it is the most important part.
  \item \textbf{Write every substitution in brackets.} Every one. This is the habit that carries the most marks.
  \item When you get one wrong, write it in your error log: what you did, and what you should have done.
\end{itemize}
\end{mdframed}

\vspace{4pt}
\begin{center}
\begin{tabularx}{\textwidth}{@{}l l X r@{}}
\textbf{Sitting} & \textbf{Part} & \textbf{What it drills} & \textbf{Minutes}\\[2pt]
\hline\rule{0pt}{11pt}
Sunday   & A    & Signed arithmetic, no calculator            & 25\\
Tuesday  & B, C & Substitution, and the four traps            & 30\\
Thursday & D    & ``Could that be true?''                     & 20\\
Saturday & E    & Six mixed questions, timed                  & 25\\
\hline
\end{tabularx}
\end{center}

% ------------------------------------------------ PART A
\PartHead{Part A --- Signed arithmetic}{No calculator. 25 minutes.}
\sitting{Sunday}
Work down the list. If you are unsure of a sign, draw a number line in the margin.

\begin{multicols}{2}
\begin{enumerate}[leftmargin=20pt,itemsep=5pt]
\item $7-(-1)=$ \ifANSWERS\else\underline{\hspace{1.5cm}}\fi \ans{8}
\item $-5-3=$ \ifANSWERS\else\underline{\hspace{1.5cm}}\fi \ans{-8}
\item $-4-(-9)=$ \ifANSWERS\else\underline{\hspace{1.5cm}}\fi \ans{5}
\item $6-(-6)=$ \ifANSWERS\else\underline{\hspace{1.5cm}}\fi \ans{12}
\item $(-3)(-7)=$ \ifANSWERS\else\underline{\hspace{1.5cm}}\fi \ans{21}
\item $(-2)^4=$ \ifANSWERS\else\underline{\hspace{1.5cm}}\fi \ans{16}
\item $(-2)^3=$ \ifANSWERS\else\underline{\hspace{1.5cm}}\fi \ans{-8}
\item $-2^4=$ \ifANSWERS\else\underline{\hspace{1.5cm}}\fi \ans{-16 \ (the minus is not inside the power)}
\item $-(-11)=$ \ifANSWERS\else\underline{\hspace{1.5cm}}\fi \ans{11}
\item $\dfrac{-20}{-4}=$ \ifANSWERS\else\underline{\hspace{1.5cm}}\fi \ans{5}
\item $12-(-8)+(-5)=$ \ifANSWERS\else\underline{\hspace{1.5cm}}\fi \ans{15}
\item $-(3-8)=$ \ifANSWERS\else\underline{\hspace{1.5cm}}\fi \ans{5}
\item $(-1)(-1)(-1)=$ \ifANSWERS\else\underline{\hspace{1.5cm}}\fi \ans{-1}
\item $5-2(-4)=$ \ifANSWERS\else\underline{\hspace{1.5cm}}\fi \ans{13}
\item $(-6)^2-6^2=$ \ifANSWERS\else\underline{\hspace{1.5cm}}\fi \ans{0}
\end{enumerate}
\end{multicols}

Expand and simplify.
\begin{enumerate}[leftmargin=20pt,itemsep=6pt,start=16]
\item $-3(x-4)$ \ifANSWERS\ans{$=-3x+12$}\else\ \underline{\hspace{6cm}}\fi
\item $-(2x-7)$ \ifANSWERS\ans{$=-2x+7$}\else\ \underline{\hspace{6cm}}\fi
\item $(x-5)(x+2)$ \ifANSWERS\ans{$=x^2-3x-10$}\else\ \underline{\hspace{5.4cm}}\fi
\item $(x+4)(x-6)$ \ifANSWERS\ans{$=x^2-2x-24$}\else\ \underline{\hspace{5.4cm}}\fi
\item $-2(3-x)+4$ \ifANSWERS\ans{$=2x-2$}\else\ \underline{\hspace{5.4cm}}\fi
\end{enumerate}

\ifANSWERS
\note{\textbf{Marking rule: below 17 of 20 means the sign work has not landed.} Open Session 3 with more of Part A rather than moving on. Items 6--8 and 15 are the diagnostic ones: they separate a student who knows the rule from one who has memorised a slogan.}
\fi

\newpage
% ------------------------------------------------ PART B
\PartHead{Part B --- Substitution, in brackets}{Calculator allowed. 15 minutes.}
\sitting{Tuesday}
Write the substitution in brackets \emph{before} you work anything out. Like this:
\[
f(x)=3x^2-x,\qquad f(-2)=3(-2)^2-(-2)=3(4)+2=14
\]
You lost marks on the diagnostic by skipping that middle step. Do not skip it here.

\begin{multicols}{2}
\begin{enumerate}[leftmargin=20pt,itemsep=7pt]
\item $f(x)=3x^2-x$. \ Find $f(-1)$. \ans{$=3(1)+1=4$}
\item $g(x)=x^2-5x$. \ Find $g(-3)$. \ans{$=9+15=24$}
\item $h(x)=-x^2+4$. \ Find $h(-5)$. \ans{$=-25+4=-21$}
\item $p(x)=2x^2+3x-1$. \ Find $p(-4)$. \ans{$=32-12-1=19$}
\item $q(x)=5-x^2$. \ Find $q(-6)$. \ans{$=5-36=-31$}
\item $r(x)=x^3+2x$. \ Find $r(-2)$. \ans{$=-8-4=-12$}
\item $s(x)=-2x-x^2$. \ Find $s(-3)$. \ans{$=6-9=-3$}
\item $f(x)=3x^2-x$. \ Find $f(-2)$. \ans{$=14$ \ (the diagnostic question)}
\end{enumerate}
\end{multicols}

% ------------------------------------------------ PART C
\PartHead{Part C --- The four traps}{Calculator allowed. 15 minutes.}
\sitting{Tuesday}
These are the four places a sign went wrong on your diagnostic, with new numbers.

\begin{enumerate}[leftmargin=20pt,itemsep=8pt]
\item A line passes through $(3,-2)$ and $(7,6)$. What is its slope?\\
  \makebox[3.2cm][l]{A) $1$}\makebox[3.2cm][l]{B) $2$}\makebox[3.2cm][l]{C) $4$}\makebox[3.2cm][l]{D) $8$}
  \ans{\\ B. $\dfrac{6-(-2)}{7-3}=\dfrac{8}{4}=2$. A comes from reading $6-(-2)$ as $6-2$.}

\item A line passes through $(-1,5)$ and $(3,-3)$. What is its slope?\\
  \makebox[3.2cm][l]{A) $-4$}\makebox[3.2cm][l]{B) $-2$}\makebox[3.2cm][l]{C) $2$}\makebox[3.2cm][l]{D) $4$}
  \ans{\\ B. $\dfrac{-3-5}{3-(-1)}=\dfrac{-8}{4}=-2$. A comes from $3-1$ in the denominator.}

\item Which expression is equivalent to $(x+6)(x-4)$?\\
  \makebox[4.2cm][l]{A) $x^2-24$}\makebox[4.2cm][l]{B) $x^2-2x-24$}%
  \makebox[4.2cm][l]{C) $x^2+2x-24$}\makebox[4.2cm][l]{D) $x^2+10x-24$}
  \ans{\\ C. $-4x+6x=+2x$. B is the sign slip you made on diagnostic Q2.}

\item Which expression is equivalent to $(x-7)(x+3)$?\\
  \makebox[4.2cm][l]{A) $x^2-4x-21$}\makebox[4.2cm][l]{B) $x^2+4x-21$}%
  \makebox[4.2cm][l]{C) $x^2-21$}\makebox[4.2cm][l]{D) $x^2-10x-21$}
  \ans{\\ A. $3x-7x=-4x$.}

\item Which inequality gives all values of $x$ for which $7-3x\ge 19$?\\
  \makebox[3.6cm][l]{A) $x\ge-4$}\makebox[3.6cm][l]{B) $x\le-4$}%
  \makebox[3.6cm][l]{C) $x\ge4$}\makebox[3.6cm][l]{D) $x\le4$}
  \ans{\\ B. $-3x\ge12$, then divide by $-3$ \emph{and flip}. Check: $x=-5$ gives $22\ge19$, true.}

\item Which inequality gives all values of $x$ for which $4-5x<24$?\\
  \makebox[3.6cm][l]{A) $x<-4$}\makebox[3.6cm][l]{B) $x>-4$}%
  \makebox[3.6cm][l]{C) $x<4$}\makebox[3.6cm][l]{D) $x>4$}
  \ans{\\ B. $-5x<20$, divide by $-5$ and flip. Check: $x=0$ gives $4<24$, true.}

\item The graph of $y=(x-2)^2-9$ crosses the $x$-axis at two points. What are their $x$-coordinates?\\
  \makebox[4.2cm][l]{A) $-1$ and $5$}\makebox[4.2cm][l]{B) $1$ and $5$}%
  \makebox[4.2cm][l]{C) $-5$ and $1$}\makebox[4.2cm][l]{D) $2$ and $9$}
  \ans{\\ A. $(x-2)^2=9$, so $x-2=\pm3$, giving $x=5$ and $x=-1$.}

\item A line is given by $5x-2y=14$. What is the $y$-coordinate of its $y$-intercept? \spr
  \ans{\\ $-7$. Put $x=0$: $-2y=14$, so $y=-7$.}
\end{enumerate}

\newpage
% ------------------------------------------------ PART D
\PartHead{Part D --- Could that be true?}{No calculator. 20 minutes.}
\sitting{Thursday}
For each one, a student has given an answer. \textbf{Do not work the question out.}
Say whether the answer could possibly be right, and in one sentence say how you know.

This is the most valuable part of the sheet. Five of the marks you lost on the
diagnostic would have survived thirty seconds of this.

\begin{enumerate}[leftmargin=20pt,itemsep=10pt]
\item A right triangle has legs of length 9 and 12. A student says the hypotenuse is $0.013$.\\
  \underline{\hspace{\linewidth}}
  \ans{\\ No. The hypotenuse is the longest side, so it must be more than 12.}

\item Of 45 students who walk to school, 26 are in Grade 11. One walker is chosen at random.
  A student says the probability they are in Grade 11 is $1$.\\
  \underline{\hspace{\linewidth}}
  \ans{\\ No. A probability of 1 means every walker is in Grade 11, but only 26 of the 45 are.}

\item Solve $5-2x\ge11$. A student says $x\ge3$. Test $x=4$.\\
  \underline{\hspace{\linewidth}}
  \ans{\\ No. $5-8=-3$, which is not $\ge11$, so the answer refutes itself in one line.}

\item A jacket costs \$260 and is reduced by 15\%. A student says the sale price is \$299.\\
  \underline{\hspace{\linewidth}}
  \ans{\\ No. A discount makes the price smaller, so it must be under \$260.}

\item A town of 250 people grows by 8\% a year. A student says that after 3 years there are 472.\\
  \underline{\hspace{\linewidth}}
  \ans{\\ No. 8\% of 250 is about 20 a year, so three years adds roughly 60--65, not 220.}

\item A cylinder has radius 3 and height 10. A student says the volume is $30\pi$.\\
  \underline{\hspace{\linewidth}}
  \ans{\\ No. The base circle alone has area $9\pi$, and the cylinder is 10 tall, so the volume is about ten times that.}
\end{enumerate}

% ------------------------------------------------ PART E
\PartHead{Part E --- Mixed, timed}{Calculator allowed. 12 minutes, timed.}
\sitting{Saturday}
Set a timer for 12 minutes. Answer every one --- there is no penalty for a wrong answer,
so a blank is a mark thrown away. Then check each answer by substituting it back.

\begin{enumerate}[leftmargin=20pt,itemsep=8pt]
\item What value of $x$ satisfies $3(x-4)=x+2$?\\
  \makebox[3.2cm][l]{A) $5$}\makebox[3.2cm][l]{B) $7$}\makebox[3.2cm][l]{C) $-5$}\makebox[3.2cm][l]{D) $-7$}
  \ans{\\ B. $3x-12=x+2$, so $2x=14$ and $x=7$.}

\item A line passes through $(-2,3)$ and $(2,-5)$. What is its slope?\\
  \makebox[3.2cm][l]{A) $-4$}\makebox[3.2cm][l]{B) $-2$}\makebox[3.2cm][l]{C) $2$}\makebox[3.2cm][l]{D) $4$}
  \ans{\\ B. $\dfrac{-5-3}{2-(-2)}=\dfrac{-8}{4}=-2$.}

\item The function $f$ is defined by $f(x)=x^2-4x$. What is the value of $f(-3)$? \spr
  \ans{\\ $21$. \ $(-3)^2-4(-3)=9+12=21$.}

\item Which inequality gives all values of $x$ for which $8-4x>20$?\\
  \makebox[3.6cm][l]{A) $x>-3$}\makebox[3.6cm][l]{B) $x<-3$}%
  \makebox[3.6cm][l]{C) $x>3$}\makebox[3.6cm][l]{D) $x<3$}
  \ans{\\ B. $-4x>12$, divide by $-4$ and flip. Check: $x=-4$ gives $24>20$, true.}

\item Which expression is equivalent to $(x-5)(x+2)$?\\
  \makebox[4.2cm][l]{A) $x^2-10$}\makebox[4.2cm][l]{B) $x^2+3x-10$}%
  \makebox[4.2cm][l]{C) $x^2-3x-10$}\makebox[4.2cm][l]{D) $x^2-7x-10$}
  \ans{\\ C. $2x-5x=-3x$.}

\item The graph of $y=(x+1)^2-16$ crosses the $x$-axis at two points.
  What is the greater of the two $x$-coordinates? \spr
  \ans{\\ $3$. \ $(x+1)^2=16$, so $x+1=\pm4$, giving $x=3$ and $x=-5$.}
\end{enumerate}

\ifANSWERS\else
\vspace{10pt}
\PartHead{Error log}{Fill this in as you go.}
\begin{center}
\begin{tabularx}{\textwidth}{|c|c|X|X|}
\hline
\textbf{Date} & \textbf{Q} & \textbf{What I did} & \textbf{What I should have done}\\
\hline
\rule{0pt}{22pt} & & &\\ \hline
\rule{0pt}{22pt} & & &\\ \hline
\rule{0pt}{22pt} & & &\\ \hline
\rule{0pt}{22pt} & & &\\ \hline
\rule{0pt}{22pt} & & &\\ \hline
\rule{0pt}{22pt} & & &\\ \hline
\rule{0pt}{22pt} & & &\\ \hline
\end{tabularx}
\end{center}
\vspace{4pt}
\begin{center}\small\itshape
Bring this sheet and your error log to the next session.
\end{center}
\fi

\ifANSWERS
\vspace{8pt}
\PartHead{Five-minute marking routine}{Tutor}
\begin{enumerate}[leftmargin=20pt,itemsep=4pt]
\item \textbf{Score Part A only.} Below 17/20, Session 3 opens with more sign work rather than moving on.
\item \textbf{Read Part D.} One line each. It tells you whether the plausibility habit is forming, and it is the part most likely to be skipped.
\item \textbf{Pick two wrong answers from Part E} and have him narrate his thinking. Leave the rest unmarked.
\end{enumerate}
If the error log is filled in, say so out loud. He will do more of what you visibly value.
\fi

\end{document}
